\documentclass[11pt,a4paper]{article}
\usepackage[utf8]{inputenc}
\usepackage[T1]{fontenc}
\usepackage[catalan]{babel}
\usepackage{lmodern}
\usepackage[sfdefault,lining]{FiraSans}
\usepackage[a4paper,top=2.3cm,bottom=2.2cm,left=1.8cm,right=1.8cm,headheight=1.3cm,headsep=0.4cm,footskip=1cm]{geometry}
\usepackage{amsmath,amssymb}
\usepackage{icomma}
\usepackage{graphicx}
\usepackage[table]{xcolor}
\usepackage{tikz}
\usepackage{fontawesome5}
\usepackage{ccicons}
\usepackage{fancyhdr}
\usepackage{lastpage}
\usepackage{tasks}
\usepackage{needspace}
\usepackage[most]{tcolorbox}
\usepackage{array}
\usepackage[hidelinks,pdftitle={Fitxa d'exercicis de la unitat 5: Equacions. 3r d'ESO},pdfsubject={Matemàtiques 3r ESO},pdfkeywords={equacions, identitats, equacions de primer grau, denominadors, equacions de segon grau, discriminant, equacions de tercer grau, regla de Ruffini, problemes}]{hyperref}

\definecolor{navy}{HTML}{1B2748}
\definecolor{accent}{HTML}{E8603A}
\definecolor{teal}{HTML}{1F9E90}
\definecolor{soft}{HTML}{EEF1F8}
\definecolor{muted}{HTML}{5D6685}

\setlength{\parindent}{0pt}
\setlength{\parskip}{3pt}
\settasks{label=\alph*), label-format=\bfseries\color{navy}, label-width=1.3em, item-indent=1.8em, column-sep=0.6em, after-item-skip=5pt, before-skip=3pt, after-skip=4pt}

% ---------- capçalera i peu ----------
\newcommand{\llicencia}{\href{https://creativecommons.org/licenses/by-nc-sa/4.0/deed.ca}{\ccbyncsa}\ \ Llicència Creative Commons Reconeixement-NoComercial-CompartirIgual 4.0 Internacional (CC BY-NC-SA 4.0)}
\pagestyle{fancy}
\fancyhf{}
\fancyhead[L]{\raisebox{-0.35\height}{\includegraphics[height=0.95cm]{logo-gva.png}}}
\fancyhead[C]{\small\color{muted}Fitxa d'exercicis · Unitat 5: Equacions · 3r d'ESO}
\fancyhead[R]{\raisebox{-0.35\height}{\includegraphics[height=0.95cm]{logo-vermellar.png}}}
\fancyfoot[L]{\scriptsize\color{muted}\llicencia}
\fancyfoot[R]{\small\color{muted}\thepage}
\renewcommand{\headrulewidth}{0.4pt}
\renewcommand{\headrule}{\color{muted!40}\hrule height \headrulewidth}
\fancypagestyle{primera}{\fancyhf{}\renewcommand{\headrulewidth}{0pt}
  \fancyfoot[L]{\scriptsize\color{muted}\llicencia}
  \fancyfoot[R]{\small\color{muted}\thepage}}

% ---------- nivells i exercicis ----------
\newcommand{\estrelles}[1]{{\color{accent}\ifcase#1\or\faStar\faStar[regular]\faStar[regular]\or\faStar\faStar\faStar[regular]\or\faStar\faStar\faStar\fi}}
\newcounter{exer}
\newcommand{\ex}[1]{\par\needspace{4\baselineskip}\medskip\stepcounter{exer}\noindent{\color{navy}\bfseries\theexer.}\ \ {\small\estrelles{#1}}\ \ }

\newcommand{\seccio}[2]{\par\needspace{10\baselineskip}\bigskip
  \begin{tcolorbox}[colback=navy,colframe=navy,coltext=white,arc=2mm,boxrule=0pt,left=3mm,right=3mm,top=1.2mm,bottom=1.2mm]
  {\large\bfseries #1\quad #2}
  \end{tcolorbox}\vspace{-2pt}}

\newcolumntype{C}[1]{>{\centering\arraybackslash}p{#1}}

% quadre de recordatori (només a la fitxa d'exercicis)
\newcommand{\recorda}[1]{\par\smallskip\begin{tcolorbox}[colback=accent!7,colframe=accent!55,arc=2mm,boxrule=0.5pt,left=2.5mm,right=2.5mm,top=1mm,bottom=1mm]\small{\color{accent}\faLightbulb[regular]}\ \textbf{Recorda:}\ #1\end{tcolorbox}}
\newcommand{\res}[1]{#1}
\newcommand{\rest}[1]{#1}

\def\fmneed{14}
% =====================================================================
% Capçalera FentMates de les fitxes: un color per a cada curs.
% Els colors són els de la pàgina inicial (LEVELS a assets/js/materials.js);
% el logo ve de _eines/logo/pdf/ (python3 _eines/logo/genera.py).
%
% Ús, al final del preàmbul de la unitat:
%   % =====================================================================
% Capçalera FentMates de les fitxes: un color per a cada curs.
% Els colors són els de la pàgina inicial (LEVELS a assets/js/materials.js);
% el logo ve de _eines/logo/pdf/ (python3 _eines/logo/genera.py).
%
% Ús, al final del preàmbul de la unitat:
%   % =====================================================================
% Capçalera FentMates de les fitxes: un color per a cada curs.
% Els colors són els de la pàgina inicial (LEVELS a assets/js/materials.js);
% el logo ve de _eines/logo/pdf/ (python3 _eines/logo/genera.py).
%
% Ús, al final del preàmbul de la unitat:
%   \input{../comu/fentmates.tex}
%   \fmcurs{1eso}
% i, al començament del document:
%   \capcalera{Matemàtiques}{1r ESO}{Fitxa d'exercicis · Unitat 2}{Divisibilitat}
% Opcional abans de \input: \def\fmneed{14} (espai mínim abans de cada secció).
% =====================================================================
\makeatletter
\def\fm@logodir{../../logo/pdf}
\providecommand{\fmneed}{10}
\def\fm@cursnom{}

% clau del curs, color i color del text damunt del color
\newcommand{\fm@curs}[3]{\@namedef{fm@color@#1}{#2}\@namedef{fm@text@#1}{#3}}
\fm@curs{1eso}{E6B84A}{navy}
\fm@curs{2eso}{3B82F6}{white}
\fm@curs{3eso}{FF7A50}{white}
\fm@curs{4eso-a}{A78BFA}{white}
\fm@curs{4eso-b}{6D28D9}{white}
\fm@curs{1batx-mi}{9B1C1C}{white}
\fm@curs{1batx-ccss}{EF6B6B}{white}
\fm@curs{2batx-mii}{1E6B47}{white}
\fm@curs{2batx-ccss}{2BB3A3}{white}

\newcommand{\fmcurs}[1]{%
  \@ifundefined{fm@color@#1}{\PackageError{fentmates}{Curs desconegut: #1}{}}{}%
  \edef\fm@tmp{\noexpand\definecolor{curs}{HTML}{\@nameuse{fm@color@#1}}%
               \noexpand\colorlet{curstext}{\@nameuse{fm@text@#1}}}\fm@tmp
  \colorlet{cursfosc}{curs!72!black}%
  \def\fm@logo{\fm@logodir/vertical-#1.pdf}%
  \def\fm@marca{\fm@logodir/marca-#1.pdf}}

% ---------- capçalera de la primera pàgina ----------
% #1 matèria (petit), #2 curs, #3 tipus de document i unitat, #4 títol
\newcommand{\capcalera}[4]{%
  \gdef\fm@cursnom{#2}%
  \par\noindent
  \begin{tcolorbox}[colback=curs,colframe=curs,boxrule=0pt,arc=2.2mm,width=3.5cm,height=2.15cm,
      valign=center,left=3mm,right=2mm,top=1mm,bottom=1mm,nobeforeafter,box align=center]
    \color{curstext}{\footnotesize #1}\par\vspace{1pt}{\fontsize{19}{21}\selectfont\bfseries #2\par}%
  \end{tcolorbox}\hspace{5mm}%
  \begin{minipage}[c]{\dimexpr\linewidth-3.5cm-5mm-2.9cm\relax}
    {\color{cursfosc}\small\bfseries\MakeUppercase{#3}}\par\vspace{2pt}
    {\color{navy}\fontsize{20}{23}\selectfont\bfseries #4\par}
  \end{minipage}\hfill
  \begin{minipage}[c]{2.6cm}\raggedleft\includegraphics[width=2.6cm]{\fm@logo}\end{minipage}\par
  \vspace{2pt}}

% ---------- barres de secció del color del curs ----------
\renewcommand{\seccio}[2]{\par\needspace{\fmneed\baselineskip}\bigskip
  \begin{tcolorbox}[colback=curs,colframe=curs,coltext=curstext,arc=2mm,boxrule=0pt,left=3mm,right=3mm,top=1.2mm,bottom=1.2mm]
  {\large\bfseries #1\quad #2}
  \end{tcolorbox}\vspace{-2pt}}

% ---------- capçalera i peu de les pàgines ----------
% primera pàgina: logos institucionals al peu
\fancypagestyle{primera}{\fancyhf{}\renewcommand{\headrulewidth}{0pt}
  \fancyfoot[L]{\raisebox{-0.32\height}{\includegraphics[height=0.78cm]{logo-gva.png}}\hspace{5mm}%
    \raisebox{-0.32\height}{\includegraphics[height=0.78cm]{logo-vermellar.png}}}
  \fancyfoot[R]{\scriptsize\color{muted}\href{https://creativecommons.org/licenses/by-nc-sa/4.0/deed.ca}{\ccbyncsa}\ \ CC BY-NC-SA 4.0 · fentmates.org\qquad\small\thepage}}
% la resta: el símbol a l'esquerra i el curs a la dreta, amb el títol al centre com abans
\fancyhead[L]{\raisebox{-0.3\height}{\includegraphics[height=0.62cm]{\fm@marca}}}
\fancyhead[R]{\small\bfseries\color{cursfosc}\fm@cursnom}
\renewcommand{\headrule}{\color{curs!45}\hrule height \headrulewidth}
\makeatother

%   \fmcurs{1eso}
% i, al començament del document:
%   \capcalera{Matemàtiques}{1r ESO}{Fitxa d'exercicis · Unitat 2}{Divisibilitat}
% Opcional abans de \input: \def\fmneed{14} (espai mínim abans de cada secció).
% =====================================================================
\makeatletter
\def\fm@logodir{../../logo/pdf}
\providecommand{\fmneed}{10}
\def\fm@cursnom{}

% clau del curs, color i color del text damunt del color
\newcommand{\fm@curs}[3]{\@namedef{fm@color@#1}{#2}\@namedef{fm@text@#1}{#3}}
\fm@curs{1eso}{E6B84A}{navy}
\fm@curs{2eso}{3B82F6}{white}
\fm@curs{3eso}{FF7A50}{white}
\fm@curs{4eso-a}{A78BFA}{white}
\fm@curs{4eso-b}{6D28D9}{white}
\fm@curs{1batx-mi}{9B1C1C}{white}
\fm@curs{1batx-ccss}{EF6B6B}{white}
\fm@curs{2batx-mii}{1E6B47}{white}
\fm@curs{2batx-ccss}{2BB3A3}{white}

\newcommand{\fmcurs}[1]{%
  \@ifundefined{fm@color@#1}{\PackageError{fentmates}{Curs desconegut: #1}{}}{}%
  \edef\fm@tmp{\noexpand\definecolor{curs}{HTML}{\@nameuse{fm@color@#1}}%
               \noexpand\colorlet{curstext}{\@nameuse{fm@text@#1}}}\fm@tmp
  \colorlet{cursfosc}{curs!72!black}%
  \def\fm@logo{\fm@logodir/vertical-#1.pdf}%
  \def\fm@marca{\fm@logodir/marca-#1.pdf}}

% ---------- capçalera de la primera pàgina ----------
% #1 matèria (petit), #2 curs, #3 tipus de document i unitat, #4 títol
\newcommand{\capcalera}[4]{%
  \gdef\fm@cursnom{#2}%
  \par\noindent
  \begin{tcolorbox}[colback=curs,colframe=curs,boxrule=0pt,arc=2.2mm,width=3.5cm,height=2.15cm,
      valign=center,left=3mm,right=2mm,top=1mm,bottom=1mm,nobeforeafter,box align=center]
    \color{curstext}{\footnotesize #1}\par\vspace{1pt}{\fontsize{19}{21}\selectfont\bfseries #2\par}%
  \end{tcolorbox}\hspace{5mm}%
  \begin{minipage}[c]{\dimexpr\linewidth-3.5cm-5mm-2.9cm\relax}
    {\color{cursfosc}\small\bfseries\MakeUppercase{#3}}\par\vspace{2pt}
    {\color{navy}\fontsize{20}{23}\selectfont\bfseries #4\par}
  \end{minipage}\hfill
  \begin{minipage}[c]{2.6cm}\raggedleft\includegraphics[width=2.6cm]{\fm@logo}\end{minipage}\par
  \vspace{2pt}}

% ---------- barres de secció del color del curs ----------
\renewcommand{\seccio}[2]{\par\needspace{\fmneed\baselineskip}\bigskip
  \begin{tcolorbox}[colback=curs,colframe=curs,coltext=curstext,arc=2mm,boxrule=0pt,left=3mm,right=3mm,top=1.2mm,bottom=1.2mm]
  {\large\bfseries #1\quad #2}
  \end{tcolorbox}\vspace{-2pt}}

% ---------- capçalera i peu de les pàgines ----------
% primera pàgina: logos institucionals al peu
\fancypagestyle{primera}{\fancyhf{}\renewcommand{\headrulewidth}{0pt}
  \fancyfoot[L]{\raisebox{-0.32\height}{\includegraphics[height=0.78cm]{logo-gva.png}}\hspace{5mm}%
    \raisebox{-0.32\height}{\includegraphics[height=0.78cm]{logo-vermellar.png}}}
  \fancyfoot[R]{\scriptsize\color{muted}\href{https://creativecommons.org/licenses/by-nc-sa/4.0/deed.ca}{\ccbyncsa}\ \ CC BY-NC-SA 4.0 · fentmates.org\qquad\small\thepage}}
% la resta: el símbol a l'esquerra i el curs a la dreta, amb el títol al centre com abans
\fancyhead[L]{\raisebox{-0.3\height}{\includegraphics[height=0.62cm]{\fm@marca}}}
\fancyhead[R]{\small\bfseries\color{cursfosc}\fm@cursnom}
\renewcommand{\headrule}{\color{curs!45}\hrule height \headrulewidth}
\makeatother

%   \fmcurs{1eso}
% i, al començament del document:
%   \capcalera{Matemàtiques}{1r ESO}{Fitxa d'exercicis · Unitat 2}{Divisibilitat}
% Opcional abans de \input: \def\fmneed{14} (espai mínim abans de cada secció).
% =====================================================================
\makeatletter
\def\fm@logodir{../../logo/pdf}
\providecommand{\fmneed}{10}
\def\fm@cursnom{}

% clau del curs, color i color del text damunt del color
\newcommand{\fm@curs}[3]{\@namedef{fm@color@#1}{#2}\@namedef{fm@text@#1}{#3}}
\fm@curs{1eso}{E6B84A}{navy}
\fm@curs{2eso}{3B82F6}{white}
\fm@curs{3eso}{FF7A50}{white}
\fm@curs{4eso-a}{A78BFA}{white}
\fm@curs{4eso-b}{6D28D9}{white}
\fm@curs{1batx-mi}{9B1C1C}{white}
\fm@curs{1batx-ccss}{EF6B6B}{white}
\fm@curs{2batx-mii}{1E6B47}{white}
\fm@curs{2batx-ccss}{2BB3A3}{white}

\newcommand{\fmcurs}[1]{%
  \@ifundefined{fm@color@#1}{\PackageError{fentmates}{Curs desconegut: #1}{}}{}%
  \edef\fm@tmp{\noexpand\definecolor{curs}{HTML}{\@nameuse{fm@color@#1}}%
               \noexpand\colorlet{curstext}{\@nameuse{fm@text@#1}}}\fm@tmp
  \colorlet{cursfosc}{curs!72!black}%
  \def\fm@logo{\fm@logodir/vertical-#1.pdf}%
  \def\fm@marca{\fm@logodir/marca-#1.pdf}}

% ---------- capçalera de la primera pàgina ----------
% #1 matèria (petit), #2 curs, #3 tipus de document i unitat, #4 títol
\newcommand{\capcalera}[4]{%
  \gdef\fm@cursnom{#2}%
  \par\noindent
  \begin{tcolorbox}[colback=curs,colframe=curs,boxrule=0pt,arc=2.2mm,width=3.5cm,height=2.15cm,
      valign=center,left=3mm,right=2mm,top=1mm,bottom=1mm,nobeforeafter,box align=center]
    \color{curstext}{\footnotesize #1}\par\vspace{1pt}{\fontsize{19}{21}\selectfont\bfseries #2\par}%
  \end{tcolorbox}\hspace{5mm}%
  \begin{minipage}[c]{\dimexpr\linewidth-3.5cm-5mm-2.9cm\relax}
    {\color{cursfosc}\small\bfseries\MakeUppercase{#3}}\par\vspace{2pt}
    {\color{navy}\fontsize{20}{23}\selectfont\bfseries #4\par}
  \end{minipage}\hfill
  \begin{minipage}[c]{2.6cm}\raggedleft\includegraphics[width=2.6cm]{\fm@logo}\end{minipage}\par
  \vspace{2pt}}

% ---------- barres de secció del color del curs ----------
\renewcommand{\seccio}[2]{\par\needspace{\fmneed\baselineskip}\bigskip
  \begin{tcolorbox}[colback=curs,colframe=curs,coltext=curstext,arc=2mm,boxrule=0pt,left=3mm,right=3mm,top=1.2mm,bottom=1.2mm]
  {\large\bfseries #1\quad #2}
  \end{tcolorbox}\vspace{-2pt}}

% ---------- capçalera i peu de les pàgines ----------
% primera pàgina: logos institucionals al peu
\fancypagestyle{primera}{\fancyhf{}\renewcommand{\headrulewidth}{0pt}
  \fancyfoot[L]{\raisebox{-0.32\height}{\includegraphics[height=0.78cm]{logo-gva.png}}\hspace{5mm}%
    \raisebox{-0.32\height}{\includegraphics[height=0.78cm]{logo-vermellar.png}}}
  \fancyfoot[R]{\scriptsize\color{muted}\href{https://creativecommons.org/licenses/by-nc-sa/4.0/deed.ca}{\ccbyncsa}\ \ CC BY-NC-SA 4.0 · fentmates.org\qquad\small\thepage}}
% la resta: el símbol a l'esquerra i el curs a la dreta, amb el títol al centre com abans
\fancyhead[L]{\raisebox{-0.3\height}{\includegraphics[height=0.62cm]{\fm@marca}}}
\fancyhead[R]{\small\bfseries\color{cursfosc}\fm@cursnom}
\renewcommand{\headrule}{\color{curs!45}\hrule height \headrulewidth}
\makeatother

\fmcurs{3eso}

\begin{document}
\thispagestyle{primera}

\capcalera{Matemàtiques}{3r ESO}{Fitxa d'exercicis · Unitat 5}{Equacions}
\vspace{10pt}
Nom i cognoms: \rule{8.4cm}{0.4pt}\hfill Grup: \rule{1.6cm}{0.4pt}\hfill Data: \rule{2.2cm}{0.4pt}
\vspace{8pt}

\begin{tcolorbox}[colback=soft,colframe=soft,arc=2mm,left=3mm,right=3mm,top=1.5mm,bottom=1.5mm]
\small Exercicis graduats per nivells de dificultat:\quad \estrelles{1}\ bàsic \qquad \estrelles{2}\ mitjà \qquad \estrelles{3}\ avançat.\\[2pt] Resol-los al quadern mostrant tots els passos i comprova les solucions sempre que pugues. Fes els càlculs sense calculadora.
\end{tcolorbox}

% =====================================================
\seccio{1}{Equacions i identitats}

\ex{1} Indica quines igualtats són identitats i quines són equacions:\begin{tasks}(2)
  \task $3(x+2)=3x+6$
  \task $3x+2=8$
  \task $(x+1)^2=x^2+2x+1$
  \task $x^2-1=0$
  \task $5x-2x=3x$
  \task $5x-2=3x$
  \task $(x-3)(x+3)=x^2-9$
  \task $2(x-1)=2x-1$
\end{tasks}\recorda{Una \textbf{identitat} es compleix per a qualsevol valor de $x$; una \textbf{equació}, només per a alguns valors (o per a cap).}

\ex{1} Comprova si els valors indicats són solució de cada equació:\begin{tasks}(2)
  \task $3x-5=x+1$, \ $x=3$
  \task $x^2-4x=5$, \ $x=-1$ i $x=5$
  \task $2(x-1)=x+4$, \ $x=5$
  \task $x^3-2x=4$, \ $x=2$
\end{tasks}

\ex{1} En l'equació $5x^2-3x+1=2x-7$, indica el primer i el segon membre, els termes, la incògnita i el grau.

\ex{2} Indica el grau de cada equació un cop reduïda:\begin{tasks}(2)
  \task $x(x+4)=x^2-8$
  \task $(x-2)^2=x^2+4$
  \task $x^3-x(x^2+1)=2$
  \task $(x+1)(x-1)=3x$
\end{tasks}\recorda{Abans de dir el grau, opera i redueix: els termes de grau alt es poden cancel·lar.}

\ex{2} Són equivalents aquestes parelles d'equacions? Justifica-ho.\begin{tasks}(1)
  \task $2x-3=7$ \ i \ $3x+1=16$
  \task $x^2=25$ \ i \ $x=5$
\end{tasks}

\ex{2} Escriu una equació de primer grau que tinga la solució $x=-2$ i una equació de segon grau que tinga les solucions $x=3$ i $x=-3$.

\ex{3} Troba el valor de $k$ en cada cas:\begin{tasks}(1)
  \task $x=3$ és solució de $2x^2-kx+3=0$.
  \task $x=-1$ és solució de $kx^3+2x=k-6$.
\end{tasks}


% =====================================================
\seccio{2}{Equacions de primer grau}

\ex{1} Resol:\begin{tasks}(4)
  \task $x+7=12$
  \task $3x=-21$
  \task $5x-4=11$
  \task $-2x+9=1$
  \task $4x+3=2x+11$
  \task $7-3x=x-13$
  \task $6x-5=9x+7$
  \task $2x-1=5x-1$
\end{tasks}\recorda{El que suma passa restant i el que resta passa sumant; el que multiplica tot un membre passa dividint (amb el seu signe).}

\ex{1} Resol traient primer els parèntesis:\begin{tasks}(2)
  \task $3(x-2)=12$
  \task $5(x+1)=2x-7$
  \task $4-(x+3)=2x-5$
  \task $2(3x-1)=4(x+2)$
\end{tasks}

\ex{2} Resol:\begin{tasks}(2)
  \task $3(2x-1)-2(x+4)=5x-13$
  \task $5-2(3-x)=4(x+1)-11$
  \task $x-[2x-(3-x)]=7$
  \task $4(x-3)-(x-5)=2(x+1)$
\end{tasks}\recorda{Un signe $-$ davant d'un parèntesi canvia el signe de \textbf{tots} els termes de dins.}

\ex{2} Resol i digues quines equacions no tenen solució i quines en tenen infinites:\begin{tasks}(2)
  \task $3(x+2)-x=2(x+3)$
  \task $5x-2(x-1)=3x+4$
  \task $6x-2(3x+1)=-2$
  \task $7x-3(2x-1)=2x+1$
\end{tasks}

\ex{2} Resol i comprova la solució:\begin{tasks}(2)
  \task $(x+1)^{2}-(x-1)^{2}=12$
  \task $x(x+3)-x^{2}=2x+5$
\end{tasks}

\ex{3} Resol:\begin{tasks}(1)
  \task $2[3-2(x+1)]=5-3(2x-1)-x$
  \task $(x-3)(x+3)-(x-1)^{2}=2(x+1)$
\end{tasks}

\ex{3} Troba el valor de $a$ perquè l'equació $3(x-a)=2x+1$ tinga la solució $x=5$.


% =====================================================
\seccio{3}{Equacions de primer grau amb denominadors}

\ex{1} Resol:\begin{tasks}(4)
  \task $\dfrac{x}{3}=4$
  \task $\dfrac{x}{2}+\dfrac{x}{4}=6$
  \task $x-\dfrac{x}{3}=10$
  \task $\dfrac{2x}{5}-1=3$
\end{tasks}

\ex{1} Resol:\begin{tasks}(4)
  \task $\dfrac{x+2}{3}=4$
  \task $\dfrac{x-1}{2}=\dfrac{x+3}{4}$
  \task $\dfrac{3x-1}{4}=2$
  \task $\dfrac{x}{6}+1=\dfrac{x}{4}$
\end{tasks}\recorda{Multiplica \textbf{tots} els termes pel m.c.m. dels denominadors (també els que no tenen denominador) i posa els numeradors entre parèntesis.}

\ex{2} Resol:\begin{tasks}(2)
  \task $\dfrac{x+1}{2}-\dfrac{x-1}{3}=2$
  \task $\dfrac{x-2}{4}-\dfrac{x-3}{6}=\dfrac{1}{2}$
  \task $\dfrac{2x+1}{3}-\dfrac{x-2}{2}=3$
  \task $x-\dfrac{x+2}{5}=\dfrac{x+6}{2}$
\end{tasks}

\ex{2} Resol:\begin{tasks}(2)
  \task $\dfrac{x-1}{3}+\dfrac{x+2}{6}=x-2$
  \task $3-\dfrac{x-1}{4}=\dfrac{x+5}{2}$
\end{tasks}

\ex{3} Resol i interpreta el resultat:\begin{tasks}(1)
  \task $\dfrac{3(x-2)}{4}-\dfrac{2(x+1)}{3}=\dfrac{x-10}{6}$
  \task $\dfrac{x}{2}-\dfrac{x-1}{3}-\dfrac{x+1}{4}=\dfrac{1-x}{12}$
  \task $\dfrac{x+1}{2}-\dfrac{x+3}{4}=\dfrac{x}{4}-1$
\end{tasks}

\ex{3} Un alumne ha resolt l'equació $\dfrac{x-2}{3}-\dfrac{x+1}{2}=1$ així:\quad $2x-4-3x+3=6\ \Rightarrow\ -x=7\ \Rightarrow\ x=-7$. \ On s'ha equivocat? Resol-la correctament.


% =====================================================
\seccio{4}{Equacions de segon grau incompletes}

\ex{1} Escriu cada equació en la forma general $ax^2+bx+c=0$, indica els coeficients $a$, $b$ i $c$ i digues si és completa o incompleta:\begin{tasks}(2)
  \task $5x-x^2=6$
  \task $3x^2=12$
  \task $x(x-4)=0$
  \task $(x+1)^2=2x+5$
\end{tasks}

\ex{1} Resol:\begin{tasks}(3)
  \task $x^{2}-36=0$
  \task $2x^{2}-50=0$
  \task $x^{2}+9=0$
  \task $4x^{2}-9=0$
  \task $3x^{2}=27$
  \task $x^{2}-7=0$
\end{tasks}\recorda{En $ax^2+c=0$ aïllem $x^2$ i fem l'arrel: hi ha \textbf{dues} solucions, $x=\pm\sqrt{\ }$, o cap si el nombre és negatiu.}

\ex{1} Resol:\begin{tasks}(4)
  \task $x^{2}-5x=0$
  \task $3x^{2}+12x=0$
  \task $2x^{2}=7x$
  \task $-x^{2}+6x=0$
\end{tasks}\recorda{En $ax^2+bx=0$ traiem factor comú $x$: una solució sempre és $x=0$. No dividisques entre $x$!}

\ex{2} Redueix i resol:\begin{tasks}(2)
  \task $(x+3)(x-3)=7$
  \task $x(x-2)=4-2x$
  \task $(2x-1)^{2}=1$
  \task $(x+2)^{2}-4=3x$
\end{tasks}

\ex{2} Resol:\begin{tasks}(2)
  \task $\dfrac{x^{2}}{3}=12$
  \task $\dfrac{x^{2}}{4}-\dfrac{x}{2}=0$
  \task $\dfrac{x^{2}+1}{5}=2$
  \task $\dfrac{x^{2}-2}{2}-\dfrac{x^{2}+1}{3}=1$
\end{tasks}

\ex{3} Resol:\begin{tasks}(1)
  \task $(2x+3)(2x-3)-(x+1)^{2}=2(x^{2}-x)-10$
  \task $\dfrac{x(x-3)}{2}-\dfrac{x(x+1)}{3}=x$
\end{tasks}

\ex{3} La Marta resol $5x^2=3x$ dividint els dos membres entre $x$ i obté $5x=3$, és a dir, $x=\tfrac35$. Què ha fet malament? Resol-la correctament.

\ex{3} Escriu una equació de segon grau incompleta que tinga per solucions:\begin{tasks}(2)
  \task $4$ i $-4$
  \task $0$ i $3$
\end{tasks}


% =====================================================
\seccio{5}{Equacions de segon grau completes}

\ex{1} Resol amb la fórmula:\begin{tasks}(4)
  \task $x^{2}-6x+8=0$
  \task $x^{2}+2x-8=0$
  \task $x^{2}-4x-5=0$
  \task $x^{2}+7x+12=0$
\end{tasks}\recorda{$x=\dfrac{-b\pm\sqrt{b^2-4ac}}{2a}$. \ Compte amb els signes: si $b=-6$, aleshores $-b=6$ i $b^2=36$.}

\ex{1} Resol:\begin{tasks}(4)
  \task $2x^{2}-5x+2=0$
  \task $3x^{2}+x-2=0$
  \task $6x^{2}-x-1=0$
  \task $4x^{2}+4x-3=0$
\end{tasks}

\ex{1} Calcula el discriminant i digues quantes solucions té cada equació, sense resoldre-la:\begin{tasks}(4)
  \task $x^{2}-10x+25=0$
  \task $2x^{2}+3x+4=0$
  \task $x^{2}-x-1=0$
  \task $-3x^{2}+2x+1=0$
\end{tasks}

\ex{2} Resol:\begin{tasks}(4)
  \task $-x^{2}+x+6=0$
  \task $9x^{2}-6x+1=0$
  \task $x^{2}+3x+5=0$
  \task $x^{2}-2x-2=0$
\end{tasks}

\ex{2} Redueix a la forma general i resol:\begin{tasks}(2)
  \task $(x-1)(x+4)=6$
  \task $(x+2)^{2}=3x+10$
  \task $x(2x-1)=(x+1)^{2}-4$
  \task $3x(x-2)-(x-1)(x+1)=9$
\end{tasks}

\ex{2} Resol:\begin{tasks}(2)
  \task $\dfrac{x^{2}}{2}-\dfrac{x}{3}=\dfrac{4}{3}$
  \task $\dfrac{x^{2}-1}{3}+\dfrac{x}{2}=2$
\end{tasks}

\ex{3} Resol:\begin{tasks}(1)
  \task $\dfrac{(x-1)^{2}}{4}-\dfrac{x+2}{6}=\dfrac{x-3}{3}$
  \task $\dfrac{x(x+1)}{2}-\dfrac{(x-1)(x+2)}{3}=1$
\end{tasks}

\ex{3} Respon:\begin{tasks}(1)
  \task Troba $k$ perquè l'equació $x^2-4x+k=0$ tinga una única solució. Quina és?
  \task Troba $k$ perquè $x=3$ siga solució de $x^2+kx-6=0$ i calcula l'altra solució.
\end{tasks}

\ex{3} Escriu una equació de segon grau, en forma general, que tinga per solucions:\begin{tasks}(2)
  \task $2$ i $-5$
  \task $\tfrac12$ i $3$
\end{tasks}


% =====================================================
\seccio{6}{Equacions de tercer grau (Ruffini)}

\ex{1} Resol, sense desenvolupar els productes:\begin{tasks}(2)
  \task $(x-2)(x+3)(x-5)=0$
  \task $x(x+1)(2x-3)=0$
  \task $(x-4)(x^2-9)=0$
  \task $(x+1)(x^2+4)=0$
\end{tasks}\recorda{Si un producte val 0, algun dels factors ha de valdre 0.}

\ex{1} Resol traient factor comú:\begin{tasks}(2)
  \task $x^{3}-16x=0$
  \task $x^{3}+3x^{2}-10x=0$
  \task $2x^{3}-8x^{2}=0$
  \task $x^{3}+x^{2}+x=0$
\end{tasks}

\ex{2} Resol amb la regla de Ruffini:\begin{tasks}(2)
  \task $x^{3}-7x+6=0$
  \task $x^{3}-4x^{2}+x+6=0$
  \task $x^{3}+2x^{2}-5x-6=0$
  \task $x^{3}-3x^{2}-4x+12=0$
\end{tasks}\recorda{Les solucions enteres són divisors del terme independent. Quan en trobes una, $x=a$, divideix per $x-a$ amb Ruffini i resol l'equació de segon grau que queda.}

\ex{2} Resol. Atenció a les solucions dobles i a les equacions de segon grau sense solució:\begin{tasks}(2)
  \task $x^{3}-x^{2}-8x+12=0$
  \task $x^{3}+3x^{2}+3x+1=0$
  \task $x^{3}-2x^{2}+x-2=0$
  \task $x^{3}+x^{2}-2=0$
\end{tasks}

\ex{2} Passa-ho tot al primer membre i resol:\begin{tasks}(2)
  \task $x^{3}+12=4x^{2}+3x$
  \task $x(x-1)(x+1)=24$
\end{tasks}

\ex{3} Resol (alguna solució és una fracció):\begin{tasks}(2)
  \task $2x^{3}-5x^{2}-4x+3=0$
  \task $3x^{3}+4x^{2}-5x-2=0$
\end{tasks}

\ex{3} Aplica Ruffini dues vegades per a resoldre aquestes equacions de quart grau:\begin{tasks}(2)
  \task $x^{4}-10x^{2}+9=0$
  \task $x^{4}-x^{3}-7x^{2}+x+6=0$
\end{tasks}

\ex{3} Troba $k$ perquè $x=2$ siga solució de $x^3-kx^2+2x+4=0$ i, amb aquest valor, troba les altres solucions.


% =====================================================
\seccio{7}{Problemes amb equacions}

\ex{1} La suma de dos nombres consecutius és 75. Quins són?\recorda{Tria la incògnita i escriu què representa, planteja l'equació, resol-la i comprova que la solució té sentit en el problema.}

\ex{1} Si a un nombre li sumem el seu triple, obtenim 56. Quin és el nombre?

\ex{1} Una persona gasta la tercera part dels seus diners en roba i la quarta part en llibres, i encara li queden 25 €. Quants diners tenia?

\ex{1} L'Anna té 12 anys i son pare, 40. D'ací a quants anys l'edat del pare serà el doble que la de l'Anna?

\ex{2} La llargària d'un rectangle és 3 m més que el doble de l'amplària i el perímetre fa 54 m. Quines dimensions té?

\ex{2} En un corral hi ha gallines i conills. En total hi ha 40 caps i 110 potes. Quants animals hi ha de cada tipus?

\ex{2} En una guardiola hi ha 30 monedes, unes d'1 € i altres de 2 €, que sumen 46 €. Quantes monedes hi ha de cada tipus?

\ex{2} L'edat d'una mare és el triple de la del seu fill, i d'ací a 12 anys només serà el doble. Quants anys tenen ara?

\ex{2} Troba dos nombres naturals consecutius el producte dels quals siga 156.

\ex{2} Si al quadrat d'un nombre li restem el seu triple, obtenim 40. Quin pot ser el nombre?

\ex{2} Si augmentem en 3 cm el costat d'un quadrat, la seua àrea augmenta 51 cm². Quant fa el costat?

\ex{3} La hipotenusa d'un triangle rectangle fa 10 cm i un catet fa 2 cm més que l'altre. Quant fan els catets?

\ex{3} Un rectangle té 34 m de perímetre i 60 m² d'àrea. Quines dimensions té?

\ex{3} Volem 20 kg d'una mescla de cafè a 9,50 €/kg barrejant cafè de 12 €/kg amb cafè de 8 €/kg. Quants quilos de cada tipus hem de posar?

\ex{3} Dos ciclistes ixen alhora de dues ciutats separades 90 km i van l'un cap a l'altre a 20 km/h i a 25 km/h. Quant de temps tarden a trobar-se? A quina distància de la primera ciutat?

\ex{3} Una caixa de base quadrada té una altura 3 cm més gran que el costat de la base, i el seu volum és 112 cm³. Quines dimensions té?

\ex{3} El producte de tres nombres naturals consecutius és 8 vegades la seua suma. Quins són?

\vfill
\begin{center}\color{muted}\small\faCheckCircle[regular]\ \ Ànim! Si et costa algun exercici, repassa els apunts interactius de la unitat.\end{center}

\end{document}
